% Данный файл распространяется под лицензией CC BY 4.0. Текст лицензии размещён на https://creativecommons.org/licenses/by/4.0/
% (c) Кафедра системного программирования, 2025

\documentclass[a4paper]{article}

\usepackage[a4paper, top=8mm, bottom=8mm, left=8mm, right=8mm]{geometry}

\usepackage{polyglossia}
\setdefaultlanguage[babelshorthands=true]{russian}
\setotherlanguage{english}

\usepackage{fontspec}
\setmainfont{FreeSerif}
\newfontfamily{\russianfonttt}[Scale=0.7]{DejaVuSansMono}

\usepackage[tiny, compact]{titlesec}

\usepackage{titling}
\setlength{\droptitle}{-1cm}
\pretitle{\begin{center}\begin{bfseries}\Large}
\posttitle{\par\end{bfseries}\end{center}}
\preauthor{\begin{center}\normalsize}
\postauthor{\par\end{center}\vspace{-1.8cm}}

\usepackage{hyperref}
\usepackage{bookmark}
\usepackage{csquotes}

\title{Разработка асинхронных процедур Network Lock Manager для NFS Mamont}

\author{Никитин Артём Денисович}

\date{}

\begin{document}

\maketitle

\begin{flushright}
    Группа: \emph{24.Б44-мм}

    Кафедра: \emph{Системного программирования}

    Научный руководитель: \emph{Гориховский Вячеслав Игоревич}

    Консультант: \emph{Яровой Вадим Дмитриевич, Инженер по разработке ПО, YADRO}
    
    Номер семестра практики: \emph{5}
\end{flushright} 

\section{Постановка задачи}

Работа выполняется в интересах компании YADRO (TATLIN.BACKUP) в рамках проекта \textbf{NFS Mamont}.

\textbf{Целью работы является} реализация асинхронных процедур Network Lock Manager (NLMv4) для разрабатываемого NFS-сервера.

NFS Mamont создается на языке Rust как современная, безопасная и высокопроизводительная альтернатива NFS-Ganesha, функционирующая в пространстве пользователя (user space).

\section{Задачи}

\begin{itemize}
    \item Решить проблему маршрутизации callback запросов без существенных изменений архитектуры сервера.
    \item Рефакторинг реестра блокировок для учёта запросов не получивших подтверждение клиента.
    \item Добавить в существующую систему обработку следующих процедур:
    \begin{itemize}
        \item NLMPROC4\_TEST\_MSG
        \item NLMPROC4\_LOCK\_MSG
        \item NLMPROC4\_CANCEL\_MSG
        \item NLMPROC4\_UNLOCK\_MSG
        \item NLMPROC4\_GRANTED\_MSG
        \item NLMPROC4\_TEST\_RES
        \item NLMPROC4\_LOCK\_RES
        \item NLMPROC4\_CANCEL\_RES
        \item NLMPROC4\_UNLOCK\_RES
        \item NLMPROC4\_GRANTED\_RES
    \end{itemize}
\end{itemize}

\section{Описание предлагаемого решения}

\begin{itemize}
\end{itemize}

\section{Апробация}

\section{Заключение}

В ходе выполнения работы получены следующие результаты:
\begin{itemize}
\end{itemize}


Ссылки на основные пуллреквесты, иллюстрирующие выполнение работы:
\begin{itemize}
\end{itemize}

\end{document}

